% =============================================================================
% BAB IV HASIL DAN PEMBAHASAN -- teks dan contoh dari templat Sain-Tek-Kes
% =============================================================================

\chapter{HASIL DAN PEMBAHASAN (terpisah atau gabung)}

Dalam penulisan hasil dan pembahasan dapat dipisah sebagai bab Hasil dan bab
Pembahasan, atau digabung menjadi bab Hasil dan Pembahasan. Bila Hasil dan
Pembahasan disatukan dalam satu bab, data/temuan disajikan dahulu, diberi
penjelasan yang cukup untuk temuan yang penting, dilanjutkan dengan penafsiran,
kemudian dengan pembahasan.

\section{Hasil}

Hasil menampilkan data/temuan dan dapat dibagi dalam beberapa subbab sesuai
dengan tujuan. Data dapat disajikan dengan ilustrasi dalam bentuk Tabel atau
Gambar (peta, denah, foto, diagram). Ilustrasi harus mampu berdiri sendiri,
artinya mudah dipahami pembaca tanpa harus merujuk teks. Tujuannya adalah
membantu pembaca memahami data dan menarik informasi penting dengan cepat
(\textit{storytelling with data}). Semua ilustrasi harus diletakkan
sedekat-dekatnya dengan teks yang menyatakan keberadaannya. Perujukan pada
ilustrasi dinyatakan di dalam paragraf sebelum tabel atau gambar. Kata
\enquote{tabel} dan \enquote{gambar} diawali dengan huruf kapital bila diikuti
dengan nomor. Nomor tabel atau gambar berurut sesuai dengan urutan kemunculannya
dalam tubuh tulisan. Contohnya adalah sebagai berikut:

\noindent\ldots{} seperti ditunjukkan pada Gambar 5.\par
\noindent\ldots{} mendekati bentuk sigmoid (Gambar 5).\par
\noindent\ldots{} meningkat dengan pesat (Tabel 3).\par
\noindent\ldots{} (lihat Lampiran 1).\par

Perujukan pada ilustrasi yang tidak disertai dengan keterangan harus
dihindari.\par
\noindent Misalnya:\par
\noindent \enquote{Hasilnya dapat dilihat pada Tabel 1} atau
\enquote{Hasilnya disajikan pada Tabel 1}.\par
\noindent Pernyataan yang lebih baik ialah\par
\noindent \enquote{Pada Tabel 1 ditunjukkan bahwa kacang hijau lebih banyak
memancarkan spektrum biru daripada kacang tanah}.

\vspace{-4pt}
\begin{table}[htbp]
  \captionsetup{justification=centering,singlelinecheck=true}
  \caption{Judul tabel biasanya pendek tanpa diakhiri tanda titik$^{\mathrm{a}}$}
  \label{tab:contoh-struktur}
  {\fontsize{11pt}{13pt}\selectfont
  \renewcommand{\arraystretch}{1.15}%
  \begin{tabular}{@{}p{3.47cm}@{}*{4}{>{\centering\arraybackslash}p{2.63cm}@{}}}
    \toprule
    \multirow{2}{3.47cm}{\centering\textbf{Judul kolom\newline pertama}}
      & \multicolumn{2}{c}{\textbf{Judul kolom$^{\mathrm{b}}$}}
      & \multicolumn{2}{c}{\textbf{Judul kolom$^{\mathrm{b}}$}} \\
    \cmidrule(lr){2-3}\cmidrule(l){4-5}
      & \hspace{5.2pt}\textbf{Subjudul\newline kolom}
      & \hspace{5.2pt}\textbf{Subjudul\newline kolom}
      & \hspace{5.2pt}\textbf{Subjudul\newline kolom}
      & \hspace{5.2pt}\textbf{Subjudul\newline kolom} \\
    \specialrule{\lightrulewidth}{0pt}{2pt}
    & \multicolumn{4}{c}{[area informasi]} \\
    Judul baris$^{\mathrm{c}}$ & & & & \\
    \quad Subjudul baris & & & & \\
    \quad Subjudul baris & & & & \\
    Judul baris & & & & \\
    \quad Subjudul baris & & & & \\
    \quad Subjudul baris & & & & \\
    & \multicolumn{4}{c}{[area informasi]} \\
    \bottomrule
  \end{tabular}
  }
  \ipbtablenote{$^{\mathrm{a}}$[catatan kaki] [titik] Sumber [jika ada]: xxxx
    xxxx [titik] $^{\mathrm{b}}$[catatan kaki] xxxx xxxx [titik]
    $^{\mathrm{c}}$[catatan kaki] xxxx xxxx [titik]}
\end{table}

\noindent Catatan kaki diketik dengan font Times New Roman 10 pt. Antara baris pertama
catatan kaki dan garis dasar tabel diberi jarak 3 pt.

\clearpage

\begin{table}[H]
  \captionsetup{format=ipbhangonehalf,justification=justified,
    singlelinecheck=false}
  \caption[Tingkat kekerasan dan kandungan gula buah pisang ambon pada suhu
    simpan yang berbeda dan pemberian putresina]{Tingkat kekerasan dan
    kandungan gula buah pisang ambon pada suhu
    simpan yang berbeda dan pemberian putresina
    \textbf{(Jika judul tabel lebih dari satu baris, maka baris selanjutnya
    sejajar dengan kata pertama pada judul tabel)}}
  \label{tab:pisang-ambon}
  {\fontsize{11pt}{12pt}\selectfont
  \begin{tabularx}{13.63cm}{@{}Xccc@{}}
    \toprule
    \multirow{2}{*}{\textbf{Perlakuan}}
      & \multicolumn{3}{c}{\textbf{Kekerasan buah dan kandungan gula pada hari ke-}} \\
    \cmidrule(l){2-4}
      & \textbf{0} & \textbf{7} & \textbf{14} \\
    \midrule
    & \multicolumn{3}{c}{Kekerasan buah (mm 50 g$^{-1}$ detik$^{-1}$)$^{\mathrm{a}}$} \\
    Suhu simpan & & & \\
    \quad $15\,^{\circ}\mathrm{C}$ & 10,20$^{\mathrm{a}}$ & 13,40$^{\mathrm{a}}$ & 11,83$^{\mathrm{a}}$ \\
    \quad $28\,^{\circ}\mathrm{C}$ & 10,64$^{\mathrm{a}}$ & 14,22$^{\mathrm{a}}$ & 88,43$^{\mathrm{b}}$ \\
    Putresina & & & \\
    \quad Dengan putresina & 11,07$^{\mathrm{a}}$ & 13,23$^{\mathrm{a}}$ & 21,19$^{\mathrm{a}}$ \\
    \quad Tanpa putresina & 10,76$^{\mathrm{a}}$ & 14,40$^{\mathrm{a}}$ & 41,82$^{\mathrm{b}}$ \\
    \midrule
    & \multicolumn{3}{c}{Gula (\%)$^{\mathrm{a}}$} \\
    Suhu simpan & & & \\
    \quad $15\,^{\circ}\mathrm{C}$ & 0,38$^{\mathrm{a}}$ & 0,56$^{\mathrm{a}}$ & 0,73$^{\mathrm{a}}$ \\
    \quad $28\,^{\circ}\mathrm{C}$ & 0,55$^{\mathrm{a}}$ & 1,82$^{\mathrm{a}}$ & 14,41$^{\mathrm{b}}$ \\
    Putresina & & & \\
    \quad Dengan putresina & 0,53$^{\mathrm{a}}$ & 0,87$^{\mathrm{a}}$ & 6,98$^{\mathrm{a}}$ \\
    \quad Tanpa putresina & 0,40$^{\mathrm{a}}$ & 1,52$^{\mathrm{a}}$ & 6,91$^{\mathrm{a}}$ \\
    \bottomrule
  \end{tabularx}
  }
  \ipbtablenote{$^{\mathrm{a}}$Angka-angka pada kolom yang sama yang diikuti
    oleh huruf yang sama tidak berbeda nyata pada taraf uji 5\% (uji selang
    berganda Duncan).}
\end{table}

\begin{figure}[H]
  \centering
  \vspace{0.65cm}
  \fbox{\parbox[c][6.4cm][c]{0.92\textwidth}{\centering
    Sisipkan diagram batang; gunakan warna dan motif yang tetap terbaca}}
  \caption[Pemilihan warna dan motif batang pada diagram batang yang jelas
    menunjukkan keterbacaan \textit{error bar}]{Pemilihan warna dan motif
    batang pada diagram batang yang jelas
    menunjukkan keterbacaan \textit{error bar}, (a) dua batang dengan penanda
    signifikansi (*), (b) beberapa batang dengan variasi warna yang ramah buta
    warna dan motif yang berbeda namun mudah dibaca. Penanda statistika
    (asteriks atau huruf) dapat diletakkan langsung di atas balok. Notasi pada
    tiap batang sudah jelas sehingga legenda tidak dicantumkan. Legenda atau
    keterangan lain dapat dicantumkan di judul gambar. Data pada diagram adalah
    \textit{random} dan dibuat dengan Microsoft Excel for Mac Version 16.89.}
  \label{fig:diagram-batang}
\end{figure}

\clearpage

Penempatan ilustrasi (tabel dan gambar) di dalam teks memerlukan perhatian
khusus untuk memastikan kualitas visualisasi ilmiah, keterbacaan, dan estetika.
Pengaturan yang tepat menjamin pembaca dapat dengan mudah menghubungkan teks
dengan ilustrasi yang disajikan. Berikut adalah panduan umum tentang penempatan
ilustrasi.

\begin{enumerate}[label=\alph*.,leftmargin=0.75cm,labelwidth=0.5cm,
  labelsep=0.25cm,align=left,itemsep=0pt,parsep=0pt,topsep=0pt,partopsep=0pt]
  \item Kedekatan dengan teks rujukan: Setiap ilustrasi ditempatkan
    sedekat mungkin setelah kalimat di dalam teks yang merujuk atau membahas
    ilustrasi tersebut. Hal ini akan meminimalkan gangguan alur baca dan
    memudahkan pembaca mengaitkan informasi.
  \item Posisi peletakan ideal: Ilustrasi diletakkan di tengah
    halaman secara horizontal, memanfaatkan ruang yang tersedia. Ilustrasi
    diposisikan sedemikian rupa sehingga judul ilustrasi berada di satu halaman
    yang sama. Penting juga untuk tidak meninggalkan ruang kosong setelah
    ilustrasi dan memastikan ada teks/paragraf yang mengikutinya.
  \item Efisiensi ruang dan orientasi: Ilustrasi disajikan dan
    diatur agar muat dalam satu halaman. Jika diperlukan, orientasi halaman
    bisa diubah menjadi lanskap untuk mengakomodasi ilustrasi yang lebar.
    Untuk tabel, format yang pendek dan lebar lebih disukai daripada yang
    panjang dan sempit. Apabila tabel dengan baris yang banyak tidak dapat
    dihindari, judul kolom berulang (\textit{repeat header row}) bisa
    digunakan, sehingga tiap halaman sambungan memiliki judul kolom.
\end{enumerate}

\vspace{13pt}

\section{Pembahasan}

Pembahasan merupakan interpretasi atau penjelasan atas data hasil kegiatan
tugas akhir. Dalam pembahasan, atau penjelasan tersebut harus dikaitkan dengan
pustaka, terutama yang mutakhir, primer, dan kredibel. Pernyataan-pernyataan
dalam paragraf dikemas dengan baik, dimulai dari pendapat sendiri di awal
paragraf, diikuti dengan dukungan pustaka, dan diakhiri dengan kalimat
penyimpulan. Argumentasi dikemukakan secara singkat dan logis yang difokuskan
untuk menjawab tujuan dan menguji hipotesis (jika ada). Argumentasi dikemukakan
untuk menunjukkan persamaan, membahas perbedaan, dan penyebab timbulnya
perbedaan, dikaitkan dengan teori dan temuan-temuan sebelumnya yang relevan.
Pembahasan berupaya menunjukkan aspek-aspek baru yang ditemukan dan merupakan
satu kesatuan. Pembahasan juga dapat mengemukakan keterbatasan dalam kegiatan
tugas akhir yang dilaksanakan.

Terdapat tiga tingkatan kedalaman dalam pembahasan hasil penelitian, yaitu
menunjukkan kejadian, pola, dan struktur. Dalam menunjukkan kejadian, peneliti
hanya membaca fakta-fakta hasil penelitian/kajian yang diperoleh. Dalam
menunjukkan pola, peneliti memperhatikan keteraturan atau kecenderungan dari
fakta-fakta hasil penelitian/kajian. Dalam menunjukkan struktur, peneliti
menjelaskan lebih lanjut dari pola dengan menjawab pertanyaan mengapa pola itu
terjadi, dan faktor apa yang menyebabkan pola itu terjadi. Pembahasan dengan
menunjukkan struktur memerlukan dukungan teori atau rujukan hasil-hasil
penelitian/kajian terdahulu yang telah teruji secara ilmiah. Pendapat peneliti
terdahulu yang sudah diringkas dalam bab Pendahuluan atau bab Tinjauan Pustaka
tidak perlu diulang lagi, tetapi diacu saja seperlunya. Sebaiknya, dalam
pembahasan juga dapat diberikan implikasi penerapan temuan baru dan segi-segi
lain yang perlu diteliti lebih lanjut. Pembahasan diakhiri dengan kalimat
positif, tegas, dan kuat.
